

We now contrast action-space sampling with latent-space planning.
Rather than comparing absolute volumes,
we compare the probability mass assigned to feasible trajectories
under different induced distributions.

\begin{proposition}[Latent Reweighting of the Search Distribution]
\label{prop:latent_reweighting}
Let $\mathcal{M}_{\mathrm{traj}} \subset \mathcal{X}$ denote the feasible trajectory set.
Consider a latent generator
$\mathcal{W}_\theta : \mathbb{R}^{d_z} \to \mathcal{X}$
and a state-conditioned latent distribution $f_\theta(s_t)$.
Assume that the learned model approximately parameterizes feasible trajectories
in the sense that there exists $\delta \in (0,1)$ such that
\begin{equation}
P_{\mathrm{latent}}(\mathcal{M}_{\mathrm{traj}})
=
\Pr_{z \sim f_\theta(s_t)}
\big[
\mathcal{W}_\theta(z) \in \mathcal{M}_{\mathrm{traj}}
\big]
\ge 1 - \delta .
\end{equation}

Let $\Phi : \mathcal{A}^H \to \mathcal{X}$ denote the trajectory rollout map
induced by the system dynamics.
Then the ratio of feasible-trajectory probabilities satisfies
\begin{equation}
\frac{
\Pr_{z \sim f_\theta(s_t)}
\big[
\mathcal{W}_\theta(z) \in \mathcal{M}_{\mathrm{traj}}
\big]
}{
\Pr_{a_{t:t+H} \sim \mathrm{Unif}(\mathcal{A}^H)}
\big[
\Phi(a_{t:t+H}) \in \mathcal{M}_{\mathrm{traj}}
\big]
}
\ge
\exp(cH)\,(1-\delta),
\end{equation}
where $c>0$ is the constant from Lemma~\ref{lem:vanishing_mass}.
Consequently, latent planning assigns exponentially larger probability mass
to feasible trajectories relative to direct action-space sampling
as the horizon grows.
\end{proposition}

\begin{proof}
Uniform sampling over $\mathcal{A}^H$ induces a trajectory distribution
via the rollout map $\Phi$.
By Lemma~\ref{lem:vanishing_mass},
the probability of sampling an $\epsilon$-feasible trajectory satisfies
\[
\Pr_{a_{t:t+H} \sim \mathrm{Unif}(\mathcal{A}^H)}
\big[
\Phi(a_{t:t+H}) \in \mathcal{M}_{\mathrm{traj}}
\big]
\le \exp(-cH).
\]

On the other hand, latent sampling induces the pushforward distribution
$P_{\mathrm{latent}} = (\mathcal{W}_\theta)_\# f_\theta$.
By assumption,
\[
\Pr_{z \sim f_\theta(s_t)}
\big[
\mathcal{W}_\theta(z) \in \mathcal{M}_{\mathrm{traj}}
\big]
\ge 1 - \delta .
\]

Taking the ratio of the two probabilities yields the stated result.
\end{proof}

\paragraph{Discussion.}
Proposition~\ref{prop:latent_reweighting} is a conditional comparison:
given a learned latent representation that approximately captures
the feasible trajectory set,
latent planning avoids the exponential decay of probability mass
that plagues action-space search.
This advantage arises from reweighting the search distribution
toward dynamically and semantically valid regions,
rather than from exhaustive exploration of a lower-dimensional space.

Proposition~\ref{prop:latent_reweighting} establishes that latent planning
reweights the search distribution toward the feasible trajectory manifold,
thereby mitigating the exponential infeasibility of action-space exploration.
Nevertheless, feasibility alone is insufficient for long-horizon planning,
as the feasible set still contains trajectories with widely varying returns.
The following corollary formalizes the necessity of iterative latent planning.
\begin{corollary}[Necessity of Iterative Latent Planning]
\label{cor:iterative_inference}
Let $\mathcal{M}_{\mathrm{traj}} \subset \mathcal{X}$ denote the feasible trajectory set,
and let
\[
V(\tau)
=
\sum_{h=0}^{H-1} \gamma^h r(s_{t+h}, a_{t+h})
\]
denote the cumulative discounted return of a trajectory $\tau$,
where $r$ is a bounded per-step reward and $\gamma \in (0,1]$.
Let
$\mathcal{W}_\theta : \mathbb{R}^{d_z} \to \mathcal{X}$
be a learned latent trajectory generator inducing the distribution
$P_{\mathrm{latent}} = (\mathcal{W}_\theta)_\# f_\theta$.

Assume that
\[
P_{\mathrm{latent}}(\mathcal{M}_{\mathrm{traj}}) \ge 1 - \delta
\quad \text{for some } \delta \in (0,1).
\]

Then, for any fixed sample budget $N$ that grows subexponentially with
the planning horizon $H$, one-shot sampling from $P_{\mathrm{latent}}$
does not, in general, guarantee with high probability
the discovery of a near-optimal trajectory.
Iterative inference procedures that adaptively reweight
the latent distribution are therefore necessary
to concentrate probability mass on high-value feasible trajectories.
\end{corollary}
\begin{proof}
Define the optimal feasible trajectory value
\[
V^\star = \sup_{\tau \in \mathcal{M}_{\mathrm{traj}}} V(\tau).
\]
For any $\varepsilon > 0$, define the $\varepsilon$-optimal feasible set
\[
\mathcal{M}_\varepsilon
=
\{\tau \in \mathcal{M}_{\mathrm{traj}} \mid V(\tau) \ge V^\star - \varepsilon\}.
\]

By construction, $\mathcal{M}_\varepsilon$ is a measurable subset of
$\mathcal{M}_{\mathrm{traj}}$.
While the latent distribution assigns non-vanishing probability mass
to the feasible set $\mathcal{M}_{\mathrm{traj}}$,
this assumption alone does not imply any uniform lower bound on
$P_{\mathrm{latent}}(\mathcal{M}_\varepsilon)$.
In particular, the probability mass of $P_{\mathrm{latent}}$
may be dispersed across trajectories with widely varying values.

Let $\tau^{(1)}, \ldots, \tau^{(N)}$ be $N$ independent samples
drawn from $P_{\mathrm{latent}}$.
The probability of observing at least one $\varepsilon$-optimal trajectory is
\[
\Pr\Big[\exists i \le N : \tau^{(i)} \in \mathcal{M}_\varepsilon \Big]
=
1 - \big(1 - P_{\mathrm{latent}}(\mathcal{M}_\varepsilon)\big)^N.
\]

If $N$ grows subexponentially with $H$,
then unless $P_{\mathrm{latent}}(\mathcal{M}_\varepsilon)$
is bounded below by a constant independent of $H$,
the above probability cannot be guaranteed to remain bounded away from zero
as the planning horizon increases.
Therefore, without additional mechanisms that adaptively reweight
the latent distribution toward higher-value regions,
one-shot latent sampling does not provide a reliable strategy
for long-horizon optimization.
Iterative inference methods, such as importance-weighted resampling,
cross-entropy optimization, or gradient-based refinement in latent space, 
explicitly modify the sampling distribution based on observed returns,
thereby increasing the probability mass assigned to $\mathcal{M}_\varepsilon$
over successive iterations.
This establishes the necessity of iterative latent planning.
\end{proof}

